
\begin{frame}{Introduction}
\normalsize


\begin{highlightbox}[sotonmintblue!30!white]
    \textbf{Purpose of this Template:} \\
    Provide a clean, consistent, and visually aligned format for academic presentations using the University of Southampton Beamer theme.
\end{highlightbox}

\insertpause
\vspace{2mm}
\rule{\linewidth}{0.5pt}
\vspace{2mm}
\insertpause

\begin{highlightbox}[sotonneongreen!20!white]
    \textbf{Template Features:}
    \begin{smartitemize}
        \item University-approved colour scheme and branding.
        \item Highlight boxes for emphasis and clarity.
        \item Flexible layout for technical content, figures, and equations.
        \item Accommodate various imaging formats (.svg, .pdf)
        \item Simple switch to convert from print and animation
    \end{smartitemize}
\end{highlightbox}

\end{frame}



\begin{frame}{Southampton Colour Palette}
\centering

\begin{columns}[T]

    \begin{column}{0.24\textwidth}
    \begin{tikzpicture}[x=0.25cm, y=0.35cm]
    \foreach \x [count=\i] in {
        sotonuniblue,
        sotondeepblue,
        sotontealblue,
        sotonskyblue,
        sotonmintblue,
        sotongreenblue,
        sotonlimegreen,
        sotonamber,
        sotonorange,
        sotonred,
        sotonmagenta,
        sotonplum,
        sotonblue,
        sotonblueDark,
        sotonblack,
        sotongreen%
    } {
      \fill[\x] (0,-\i) rectangle +(4,0.8);
      \node[anchor=west, inner sep=2pt] at (4.1,-\i+0.45) {\texttt{\x}};
    }
    \end{tikzpicture}
    \end{column}

    \begin{column}{0.24\textwidth}
    \begin{tikzpicture}[x=0.25cm, y=0.35cm]
    \foreach \x [count=\i] in {
        sotongreen,
        sotonlightgray,
        sotonanthrazit,
        sotonlink,
        sotonaqua,
        sotonlightpurple,
        sotondarkpurple,
        sotonpurple,
        sotonlilac,
        sotonblueceleste,
        sotondarkblue,
        sotonlightblue,
        sotonbrightblue,
        sotonnavyblue,
        sotonneongreen,
        sotondarkgreen%
    } {
      \fill[\x] (0,-\i) rectangle +(4,0.8);
      \node[anchor=west, inner sep=2pt] at (4.1,-\i+0.45) {\texttt{\x}};
    }
    \end{tikzpicture}
    \end{column}

    \begin{column}{0.24\textwidth}
    \begin{tikzpicture}[x=0.25cm, y=0.35cm]
    \foreach \x [count=\i] in {
        sotonmidgreen,
        sotonbrightgreen,
        sotonlightgreen,
        sotontruegreen,
        sotonstone,
        sotondarkgrey,
        sotondarkbrown,
        sotondarkred,
        sotonburgundy,
        sotonpink,
        sotonrichred,
        sotonmidred,
        sotonyellow,
        sotonferrarired,
        sotonneonyellow,
        sotontruered%
    } {
      \fill[\x] (0,-\i) rectangle +(4,0.8);
      \node[anchor=west, inner sep=2pt] at (4.1,-\i+0.45) {\texttt{\x}};
    }
    \end{tikzpicture}
    \end{column}

    \begin{column}{0.24\textwidth}
    \begin{tikzpicture}[x=0.25cm, y=0.35cm]
    \foreach \x [count=\i] in {
        sotontrueblue,
        sotontrueyellow,
        sotonplainblack,
        sotoninkblack,
        sotonslategrey,
        sotonsteelgrey,
        sotoncloudgrey,
        sotonmistgrey,
        sotonneonpink,
        sotonneonorange,
        sotonneonblue,
        sotonneonpurple,
        sotonneonred%
    } {
      \fill[\x] (0,-\i) rectangle +(4,0.8);
      \node[anchor=west, inner sep=2pt] at (4.1,-\i+0.45) {\texttt{\x}};
    }
    \end{tikzpicture}
    \end{column}

\end{columns}
\end{frame}




\begin{frame}{Example with Two Columns}
\small 

Demonstration of using two columns to present related content side-by-side.

\vspace{2mm}
\begin{columns}[T] % [T] aligns the top of the columns

    \begin{column}{0.48\textwidth}
    \begin{highlightbox}[sotonneongreen!30!white]
        \textbf{Text + Image Example}

        You can place a description, bullet points, or equations here to complement the visual on the right.
    \end{highlightbox}

    \vspace{2mm}
    \centering
    \includegraphics[width=0.6\linewidth]{example-image-a} % Replace with actual image
    \end{column}

    \begin{column}{0.48\textwidth}
    \begin{highlightbox}[sotonuniblue!30!white]
        \textbf{Visual Emphasis}

        Use visuals to present results, schematics, or concepts clearly.
    \end{highlightbox}

    \vspace{2mm}
    \centering
    \includegraphics[width=0.6\linewidth]{example-image-b} % Replace with actual image
    \end{column}

\end{columns}
\end{frame}

\begin{frame}{Example Slide}
\footnotesize

This slide demonstrates how to present example numerical data using tables and highlight boxes in a consistent layout.

\insertpause

\begin{table}[H]
\renewcommand{\arraystretch}{1.4}
\centering
\begin{tabular}{|l|c|c|c|c|c|c|}
\hline
Label & A & B (units) & C (units) & D (units) & E (units) & Category \\
\hline
\multirow{2}{*}{X1} & \multirow{2}{*}{10} & \multirow{2}{*}{123.4} & \multirow{2}{*}{5.67} & \multirow{2}{*}{0.0089} & 2500.00 & Type A \\
                    &                     &                        &                       &                          & 2480.50 & Type B \\
\hline
\multirow{2}{*}{X2} & \multirow{2}{*}{12} & \multirow{2}{*}{456.7} & \multirow{2}{*}{3.21} & \multirow{2}{*}{0.0065} & 3100.00 & Type A \\
                    &                     &                        &                       &                          & 3088.00 & Type B \\
\hline
\end{tabular}
\end{table}

\insertpause
\begin{highlightbox}[sotonmintblue!30!white]
\centering
\textbf{Example Highlight:} Value E changes across categories
\end{highlightbox}

\insertpause
\begin{highlightbox}[sotonneongreen!30!white]
\centering
\textbf{Example Highlight:} Label X2 corresponds to higher A and lower D
\end{highlightbox}

\insertpause
\begin{highlightbox}[sotonamber!30!white]
\centering
\textbf{Example Highlight:} Category B generally yields lower E values \citep{Lovell:2011:updated}
\end{highlightbox}

\tiny
\vfill
\vspace{-2mm}
\textsuperscript{1} Lovell, C. J. (2011). Updated templates

\end{frame}